\documentclass[11pt]{article}
\usepackage[margin=2.5cm]{geometry}
\emergencystretch=3em
\usepackage{amsmath,amssymb}
\usepackage{amsthm}
\usepackage{booktabs}
\usepackage{graphicx}
\usepackage{float}
\usepackage{hyperref}
\hypersetup{colorlinks=true, linkcolor=black, citecolor=black, urlcolor=black}

\newtheorem{proposition}{Proposition}
\newtheorem{definition}{Definition}
\newtheorem{remark}{Remark}
\renewcommand{\proofname}{Proof}

\title{\Large If an ETF holds illiquid securities, how can the ETF itself be liquid?\\[4pt]
\large A statistical test of liquidity transformation via NAV smoothing\\[2pt]
\normalsize Satellite analysis -- LQD and HYG}
\author{Tommaso Zeri \\ \small Advanced Fixed Income and Credit --- Supplementary analysis \\ \small Theme: \emph{Market microstructure insights}}
\date{September 2026}

\begin{document}
\maketitle

\begin{abstract}
\noindent A bond ETF holds hundreds of corporate securities that, individually, trade infrequently and at irregular intervals. Yet the ETF itself trades continuously on exchange, with tight spreads. This note builds and empirically verifies a direct statistical test of this phenomenon: we formalize two reduced-form models -- one for the Net Asset Value (NAV), one for the ETF's exchange price -- that translate into testable hypotheses the idea that the NAV inherits the illiquidity of the underlying basket while the price, forming continuously on the secondary market, does not. The gap between the statistical properties of the two processes, measured via serial autocorrelation and formally tested with the Ljung--Box test, constitutes the empirical evidence that answers the research question. We validate the estimators on synthetic data with a known staleness parameter, then apply them to the real historical NAV of LQD (investment-grade) and HYG (high-yield) and, where OHLCV data are available, to the direct NAV-versus-price comparison for both funds.
\end{abstract}

% ============================================================================
\section{Research question and framing}
\label{sec:rq}

\textbf{Research question.} If an ETF holds an illiquid underlying, how can the ETF itself be liquid?

This note is a satellite analysis, complementary to the main analysis that answers the question by laying it out across the ETF's two market levels: the \emph{primary market}, where Authorized Participants (APs) create and redeem shares in blocks (\emph{creation units}) through arbitrage against the NAV; and the \emph{secondary market}, order-driven, where the ETF's price forms continuously thanks to market making, with the market maker managing its own inventory risk through \emph{continuous hedging} on correlated liquid proxies, rather than directly and continuously trading the individual illiquid bonds underneath.

This mechanistic account -- \emph{how} the ETF stays liquid -- is developed in the main analysis. This note addresses a complementary, logically prior question: \emph{can we measure, in the data, that this transformation is actually taking place?} The answer we build here is statistical: we show that the NAV carries a quantifiable fingerprint of the underlying's illiquidity (positive serial autocorrelation of returns) that the ETF's exchange price does not share, and that this gap is statistically significant, not a sampling artifact. The \emph{why} behind this gap -- the AP and market-maker mechanism -- remains the content of the main analysis; this note puts a number on it.

% ============================================================================
\section{Methodology: the two models}
\label{sec:models}

\subsection{Model 1: NAV smoothing (Getmansky, Lo and Makarov, 2004)}
\label{sec:model1}

The idea, originally developed for illiquid hedge funds and private equity, applies unchanged to a bond ETF. Let $r^\ast_t$ be the ``true'' (unobservable) return of the bond basket on day $t$: the return that would be observed if every security in the basket repriced instantaneously on the basis of all available information. The return \textbf{reported} in the NAV is not $r^\ast_t$, but a weighted average of recent true returns,
\[
r^{NAV}_t = \sum_{j=0}^{k-1} \theta_j\, r^\ast_{t-j}, \qquad \theta_j \ge 0,\quad \sum_{j=0}^{k-1}\theta_j = 1,
\]
because when a bond does not trade today, the pricing desk uses a matrix price based on the last observed price -- stale information that ``leaks'' into today's NAV.

\paragraph{Why this creates autocorrelation: a minimal example.} The simplest possible case: half the weight on today, half on yesterday,
\[
r^{NAV}_t = 0.5\, r^\ast_t + 0.5\, r^\ast_{t-1}.
\]
With (independent) true returns on days 1--4 equal to $+2\%,\ -1\%,\ +3\%,\ +1\%$:
\begin{align*}
r^{NAV}_2 &= 0.5(-1\%)+0.5(2\%) = 0.5\%, \\
r^{NAV}_3 &= 0.5(3\%)+0.5(-1\%) = 1.0\%, \\
r^{NAV}_4 &= 0.5(1\%)+0.5(3\%) = 2.0\%.
\end{align*}
$r^{NAV}_3$ and $r^{NAV}_4$ share the same ingredient -- the $+3\%$ shock from day 3, once as ``today'' in $r^{NAV}_3$, again as ``yesterday'' in $r^{NAV}_4$. That shock is not consumed in one shot: half of it lands in today's NAV, the other half also drags tomorrow's NAV upward. It is precisely this overlap of ingredients between consecutive observations that generates positive autocorrelation. Without smoothing ($r^{NAV}_t = r^\ast_t$), there is no shared ingredient between one day and the next, and the autocorrelation is zero.

\paragraph{General parametric case.} Geometrically decaying weights,
\[
\theta_j = (1-\phi)\phi^{\,j}, \qquad \phi \in [0,1),
\]
where $\phi$ is the \textbf{smoothing parameter} ($\phi=0$: no smoothing, clean NAV; $\phi \to 1$: very persistent smoothing, very stale NAV). Assume $r^\ast_t$ is i.i.d.\ with mean zero and variance $\sigma^2$.

\begin{proposition}[Lag-1 autocorrelation of the NAV]
\label{prop:rho1}
In the limit $k\to\infty$, $r^{NAV}_t$ is an $MA(\infty)$ process with autocovariance
\[
\gamma(h) = \sigma^2\,\frac{(1-\phi)^2}{1-\phi^2}\,\phi^{\,h}, \qquad h\ge 0,
\]
so that $\rho_1 := \gamma(1)/\gamma(0) = \phi$ \emph{exactly}, and $\gamma(0) = \sigma^2\frac{1-\phi}{1+\phi} < \sigma^2$.
\end{proposition}

\begin{proof}
By definition of autocovariance and with $r^\ast_t$ i.i.d.\ (so $\mathrm{Cov}(r^\ast_{t-i}, r^\ast_{t-j}) = \sigma^2 \mathbb{1}\{i=j\}$),
\[
\gamma(h) = \mathrm{Cov}(r^{NAV}_t,\, r^{NAV}_{t-h}) = \mathrm{Cov}\Big(\sum_{i\ge 0}\theta_i r^\ast_{t-i},\ \sum_{j\ge 0}\theta_j r^\ast_{t-h-j}\Big) = \sum_{i,j\ge 0}\theta_i\theta_j\,\mathrm{Cov}(r^\ast_{t-i}, r^\ast_{t-h-j}).
\]
The covariance term is nonzero only when $t-i = t-h-j$, i.e.\ $i = h+j$; the double sum collapses to
\[
\gamma(h) = \sigma^2 \sum_{j\ge 0} \theta_{h+j}\,\theta_j.
\]
Substituting $\theta_j = (1-\phi)\phi^{\,j}$:
\[
\gamma(h) = \sigma^2(1-\phi)^2\phi^{\,h}\sum_{j\ge 0}\phi^{2j} = \sigma^2(1-\phi)^2\phi^{\,h}\cdot\frac{1}{1-\phi^2},
\]
using the geometric-series identity $\sum_{j\ge 0}\phi^{2j} = 1/(1-\phi^2)$ for $|\phi|<1$. This gives $\gamma(h) = \sigma^2\frac{(1-\phi)^2}{1-\phi^2}\phi^{\,h}$. For $h=0$: $\gamma(0) = \sigma^2\frac{(1-\phi)^2}{1-\phi^2} = \sigma^2\frac{1-\phi}{1+\phi}$, smaller than $\sigma^2$ for every $\phi>0$ -- smoothing \emph{dampens} observed volatility, in addition to introducing autocorrelation (Getmansky-Lo-Makarov's original empirical point: stale returns also produce inflated Sharpe ratios). The ratio $\gamma(1)/\gamma(0)$ simplifies immediately to $\phi$.
\end{proof}

\begin{remark}[Why $\hat\phi = \hat\rho_1$ is a valid Yule--Walker estimator]
The Yule--Walker equations for an $AR(1)$ process give $\phi = \rho_1$ as a consequence of its recursive structure ($r_t = \phi r_{t-1} + \varepsilon_t$). Our $r^{NAV}_t$ is instead an $MA(\infty)$, generated by aggregating past shocks, not by autoregressive feedback -- yet Proposition~\ref{prop:rho1} shows that its autocovariance function $\gamma(h)\propto\phi^h$ is \emph{identical} to that of an $AR(1)$ with the same $\phi$. The estimator $\hat\phi=\hat\rho_1$ is therefore legitimate not because the NAV is generated by an $AR(1)$, but because it shares the same autocovariance fingerprint -- a modeling choice that is more economically faithful (the mechanism is a smoothing of past prices, not feedback in the return itself) while being statistically indistinguishable in this lag-1 diagnostic.
\end{remark}

\subsection{Model 2: the ETF price on the secondary market}
\label{sec:model2}

The ETF's exchange price is modeled as
\[
r^{P}_t = r^\ast_t + \eta_t, \qquad \eta_t \sim \text{i.i.d.}(0,\sigma_\eta^2) \text{ independent of } r^\ast_t,
\]
where $\eta_t$ represents the idiosyncratic noise of a price that forms continuously through order flow on the secondary market. Being the sum of two independent i.i.d.\ processes, $r^P_t$ is itself i.i.d., so $\rho_1(r^P) = 0$ in population: no day-on-day statistical memory, in sharp contrast to Model 1.

\begin{remark}[An honest idealization]
A market with a bid-ask spread typically generates mild \emph{negative} autocorrelation in low-frequency returns (Roll, 1984, bid-ask bounce), not an exactly zero $\rho_1$. Model 2 is therefore an idealization: it predicts $\rho_1(P)\approx 0$, not $\rho_1(P)=0$ to the decimal. As shown in Section~\ref{sec:results-gap}, the real data are consistent with this idealization up to this well-known deviation -- which, if anything, strengthens rather than weakens the gap relative to the NAV.
\end{remark}

\subsection{Estimators and tests}
\label{sec:estimators}

\paragraph{Sample lag-1 autocorrelation.} Over a rolling window of $W=60$ trading days, $\hat\rho_1 = \widehat{\mathrm{Corr}}(r_t, r_{t-1})$, recomputed on a rolling basis along the series.

\paragraph{Variance ratio (Lo and MacKinlay, 1988).}
\[
VR(q) = \frac{\mathrm{Var}\big(\sum_{i=0}^{q-1} r_{t-i}\big)}{q\cdot\mathrm{Var}(r_t)},
\]
equal to $1$ under i.i.d.\ returns and increasing in $\phi$ for the geometric process of Model 1 (the autocovariance decays as $\phi^h$, exactly as for an $AR(1)$). We use $q=5$.

\paragraph{Ljung--Box test.} To formally test $H_0: \rho_1=\cdots=\rho_L=0$ (no autocorrelation up to lag $L$) against $H_1:$ at least one $\rho_k\ne 0$, the statistic
\[
Q = n(n+2)\sum_{k=1}^{L}\frac{\hat\rho_k^2}{n-k}
\]
is asymptotically $\chi^2_L$ under $H_0$. We use $L=10$. Unlike $\hat\rho_1$ alone, the Ljung--Box test aggregates information across multiple lags and delivers a formal $p$-value: it tells us whether the observed autocorrelation is too strong to be explained by sampling noise alone, not just how large it is on average.

% ============================================================================
\section{How the analysis formally answers the research question}
\label{sec:formal-answer}

The two models are not a self-contained statistical exercise: they are the translation into testable hypotheses of the two sides of the liquidity-transformation mechanism, and their joint verification constitutes the formal answer to the research question.

\begin{definition}[Liquidity-transformation gap]
For a given fund, define
\[
\widehat{\Delta} \;=\; \hat\rho_1(NAV) - \hat\rho_1(P).
\]
Under Model 1 and Model 2, $\widehat{\Delta}\xrightarrow{p}\phi - 0 = \phi$: the gap converges to the underlying's staleness parameter.
\end{definition}

The logic of the test proceeds in three steps.

\begin{enumerate}
\item \textbf{Is the underlying illiquid?} We test $H_0:$ no autocorrelation in the NAV, via Ljung--Box. If we reject $H_0$ with strong significance, we are not assuming the underlying is illiquid -- we are concluding it from the data: only pricing based on stale information produces that autocorrelation pattern.
\item \textbf{Is the ETF price ``clean''?} We test the same hypothesis on the ETF's exchange price. If we fail to reject $H_0$ (or the estimated effect is economically negligible and/or of the opposite sign, consistent with bid-ask bounce), the ETF price behaves like a process close to i.i.d. -- no statistical inheritance from the underlying's smoothing.
\item \textbf{Is the gap real?} If point (1) holds and point (2) holds simultaneously, for the same fund, over the same period, then $\widehat{\Delta} > 0$ cannot be explained simply by saying ``the ETF and its underlying are the same asset'' -- if they were, they would inherit the same autocorrelation. The gap therefore measures, with a number rather than only a narrative, \emph{how much} the ETF wrapper is decoupling the liquidity observed on the secondary market from the actual liquidity of the underlying.
\end{enumerate}

The outcome of this logical chain (Section~\ref{sec:results}) answers \emph{yes} to the research question, with statistical evidence: an ETF can be liquid even while holding an illiquid underlying because the two price levels -- valuation NAV and exchange price -- do not share the same generating process, and this non-sharing is measurable and statistically robust, not merely plausible at a narrative level. The \emph{how} -- the market maker's continuous hedging on liquid proxies, the AP's creation/redemption periodically anchoring $P_t$ to $NAV_t$ without requiring every trade to touch the underlying -- is the economic mechanism developed in the main analysis; this note supplies its quantitative counterpart.

% ============================================================================
\section{Results}
\label{sec:results}

\subsection{Validation on synthetic data}
\label{sec:results-synth}

We simulate i.i.d.\ true returns ($n=20{,}000$), smooth them with a known $\phi$ to obtain the synthetic NAV, and generate a synthetic ETF price as the true return plus independent idiosyncratic noise. Table~\ref{tab:validation} shows that the NAV's sample autocorrelation recovers $\phi$ almost exactly, with Ljung--Box significance rising sharply with $\phi$; the ETF price's autocorrelation stays close to zero and is never significant, by construction.

\begin{table}[H]
\centering
\caption{Recovery of $\phi$ on long synthetic samples ($n=20{,}000$, $L=10$ Ljung--Box lags).}
\label{tab:validation}
\begin{tabular}{ccccc}
\toprule
$\phi$ (true) & $\hat\rho_1(NAV)$ & Ljung--Box $p$ (NAV) & $\hat\rho_1(ETF)$ & Ljung--Box $p$ (ETF) \\
\midrule
$0.0$ & $-0.0054$ & $0.5901$ & $-0.0047$ & $0.5799$ \\
$0.2$ & $\phantom{-}0.1957$ & $<0.0001$ & $-0.0062$ & $0.4759$ \\
$0.4$ & $\phantom{-}0.3876$ & $<0.0001$ & $-0.0132$ & $0.3616$ \\
$0.6$ & $\phantom{-}0.5986$ & $<0.0001$ & $-0.0031$ & $0.9267$ \\
$0.8$ & $\phantom{-}0.8052$ & $<0.0001$ & $\phantom{-}0.0181$ & $0.3368$ \\
\bottomrule
\end{tabular}
\end{table}

At $\phi=0$ (no smoothing), Ljung--Box correctly \emph{fails to reject} $H_0$ for the NAV ($p=0.5901$): the estimator finds no autocorrelation where there is none. For every $\phi>0$ tested, it rejects $H_0$ with $p<0.0001$. On the ETF price, $H_0$ is never rejected, regardless of $\phi$: the gap does not depend on how stale the NAV is, only its NAV side does.

A second test uses a synthetic path with a time-varying smoothing regime ($\phi$ low $\to$ high $\to$ low $\to$ high $\to$ low, across 5 segments of 500 days), to verify that the \emph{rolling} estimates (60-day window) track the regime shift without any explicit classification step.

\begin{figure}[H]
\centering
\includegraphics[width=0.95\textwidth]{../figures/smoothing_rolling_diagnostics_en.png}
\caption{Top panel: known smoothing regime (ground truth). Middle panel: rolling $\hat\rho_1$ of the NAV (red) visibly tracks the known regime, while that of the ETF price (blue) stays close to zero regardless of the regime. Bottom panel: the rolling variance ratio shows the same pattern.}
\label{fig:rolling}
\end{figure}

\textbf{A note on interpretation.} A large gap is \emph{not} a warning sign: it means the underlying is stale \emph{and that, despite this, the ETF price stays clean} -- the liquidity transformation is working as intended. A genuine breakdown of the mechanism would look different: not $\hat\rho_1(NAV)$ rising, but $\hat\rho_1(ETF)$ itself starting to rise toward the NAV's level -- the gap would close ``from the wrong side.'' This is precisely the pattern we observe in real data during crises (Section~\ref{sec:results-gap}).

\subsection{Real data: NAV staleness, LQD versus HYG}
\label{sec:results-nav}

Using the real historical NAV (data already exported from iShares), we compute $\hat\rho_1$, $VR(5)$, and the Ljung--Box $p$-value on the NAV side for both funds. The economic hypothesis: high-yield bonds are typically less liquid than investment-grade bonds, so HYG's NAV should show \emph{more} smoothing.

\begin{table}[H]
\centering
\caption{NAV staleness diagnostics, real data ($L=10$ Ljung--Box lags).}
\label{tab:real-nav}
\small
\begin{tabular}{lccccc}
\toprule
Fund & N days & $\hat\rho_1(NAV)$ & $VR(5)$ & $\hat\rho_1(NAV)$ & Ljung--Box \\
 & & full-sample & full-sample & rolling mean & $p$ \\
\midrule
LQD (investment-grade) & 6{,}079 & $0.0999$ & $1.2201$ & $0.0235$ & $<0.0001$ \\
HYG (high-yield) & 4{,}897 & $0.3409$ & $1.9647$ & $0.2126$ & $<0.0001$ \\
\bottomrule
\end{tabular}
\end{table}

\begin{figure}[H]
\centering
\includegraphics[width=0.95\textwidth]{../figures/smoothing_real_nav_rho1_en.png}
\caption{Rolling (60-day) NAV autocorrelation, LQD (top) and HYG (bottom). HYG stays visibly higher on average.}
\label{fig:real-nav}
\end{figure}

The result confirms the hypothesis: HYG's $\hat\rho_1(NAV)$ and $VR(5)$ are both more than three times LQD's, and both highly significant -- direct evidence, based on NAV data alone, that the high-yield basket is structurally more stale/illiquid than the investment-grade one.

\subsection{Real data: direct comparison of NAV against ETF price}
\label{sec:results-gap}

Repeating the same diagnostic on the ETF's exchange price (OHLCV) and comparing it directly with the NAV, for the same fund over the same period, gives the most direct test of the research question: the gap $\widehat\Delta$ defined in Section~\ref{sec:formal-answer}.

\begin{table}[H]
\centering
\caption{NAV-versus-price comparison, real data (means of the 60-day rolling estimates).}
\label{tab:real-gap}
\begin{tabular}{lccc}
\toprule
Fund & $\hat\rho_1(NAV)$, mean & $\hat\rho_1(P)$, mean & $\widehat\Delta$, mean \\
\midrule
LQD (investment-grade) & $0.0233$ & $-0.0454$ & $0.0687$ \\
HYG (high-yield) & $0.2124$ & $-0.0051$ & $0.2175$ \\
\bottomrule
\end{tabular}
\end{table}

Three observations. First, $\widehat\Delta$ is positive for both funds and roughly three times larger for HYG than for LQD -- consistent with Table~\ref{tab:real-nav} and with the IG-versus-HY hypothesis, this time confirmed on the price side as well, not only the NAV. Second, $\hat\rho_1(P)$ is not exactly zero but slightly \emph{negative} for both funds: consistent with the honest idealization discussed in Section~\ref{sec:model2} (Roll's bid-ask bounce), and if anything this strengthens the gap rather than weakening it, since NAV and price not only fail to share autocorrelation, they move in opposite directions.

Third, and more interesting: looking at the \emph{rolling} path of $\hat\rho_1(NAV)$ and $\hat\rho_1(P)$ over time (not shown in the table), the gap is not constant. During the systemic stress episodes of 2008 and 2020, $\hat\rho_1(P)$ rises together with $\hat\rho_1(NAV)$, instead of staying close to zero -- exactly the ``genuine breakdown'' pattern anticipated in Section~\ref{sec:results-synth}: under stress, the market maker struggles to maintain low-cost continuous hedging (wider spreads, lower willingness to hold inventory), and some of the underlying's staleness leaks into the ETF price itself. Liquidity transformation is therefore not a constant property of the ETF wrapper, but is \emph{state-dependent}: stronger in normal conditions, weaker precisely when it would be needed most.

% ============================================================================
\section{Discussion}
\label{sec:discussion}

Three independent lines of evidence converge on the same answer. On synthetic data (Section~\ref{sec:results-synth}), the estimator correctly recovers the known staleness parameter and the Ljung--Box test declares it significant exactly when, and only when, smoothing is present -- the measurement instrument works. On real NAV-only data (Section~\ref{sec:results-nav}), the a priori expected IG$<$HY liquidity gradient emerges with strong statistical significance. On the direct NAV-versus-price comparison (Section~\ref{sec:results-gap}), the same gradient reappears in the gap $\widehat\Delta$, with the sign and magnitude predicted by Model 2, and with a time dynamic (weakening during crises) consistent with the economic mechanism -- not postulated in the models, but observable in the data.

The answer to the research question is therefore not merely narrative: a bond ETF can be liquid despite holding an illiquid underlying because its exchange price and its NAV are, statistically, two different processes -- and this difference is measurable, statistically significant, and consistent in direction, magnitude, and time dynamics with the primary-market/secondary-market, market-making, and creation/redemption mechanism developed in the main analysis.

% ============================================================================
\section{Limitations}
\label{sec:limits}

The rolling window (60 days) is a design choice subject to the usual bias/variance trade-off; the geometric smoothing model is a convenient parametrization, not the only possible one; Model 2 is an idealization that neglects the bid-ask bounce (Roll, 1984), consistent with the slight negative autocorrelation observed in the real data; a positive $\hat\rho_1(NAV)$ could in principle reflect sources of autocorrelation other than smoothing, not isolated in this note; the direct comparison in Section~\ref{sec:results-gap} requires the ETF's exchange OHLCV (network access), not always available in every execution environment -- Section~\ref{sec:results-nav} (NAV only) remains a self-contained real-data test of the IG-versus-HY hypothesis regardless.

% ============================================================================
\section{Conclusions}
\label{sec:conclusions}

A bond ETF's NAV inherits the illiquidity of the underlying bonds in the form of positive serial autocorrelation, statistically significant (Ljung--Box, $p<0.0001$ on both LQD and HYG). The exchange price, forming through continuous trading among investors on the secondary market, does not inherit the same property -- if anything it shows the mild negative autocorrelation expected of a quoted market with a bid-ask spread. The gap between the two, positive and consistent in magnitude with the lower liquidity of high-yield credit, is the direct empirical evidence that an ETF can be liquid in the sense of having a clean, continuously updated price, even though the value that price tracks moves in a structurally slower way -- and that this transformation, while robust under normal conditions, weakens precisely during episodes of systemic stress in which the hedging and arbitrage mechanism that sustains it is put under pressure.

\end{document}
